\section{One run, start to finish}

Every formula in this guide is small. What is easy to lose is the chain: a balance reading
becomes litres, litres and a lab concentration become milligrams, two sets of milligrams
become a crossing, and only then does anything look like a result. This section carries run
AE through that chain with the numbers the code actually produces, so any step can be checked
on paper. Run AG follows in a table beside it.

Every figure below came from running the code, not from reading the formulas. A reader who
wants to repeat it can print any intermediate value the same way:

\begin{quote}
\texttt{uv run python -c "... AmountThatCrossed} \\
\texttt{.find\_milligrams(run, CONCENTRATE, LITHIUM)"}
\end{quote}

\subsection{What AE arrived with}

AE used a GO-10 cation membrane over $6.0\unit{cm^2}$, fed at $400\unit{mg/L}$ lithium and
$2000\unit{mg/L}$ magnesium, and ran for 240 minutes. Only the readings in
Table~\ref{tab:ae-inputs} reach the two outcomes the model is fitted to. Everything else in
AE's eighty-odd rows is there for the checks, for the plots, or for nothing at all.

\begin{table}[htbp]
\centering
\caption{The readings AE's result is built from. The tares come from the run sheet; the
weights from the balance under each tank; the concentrations from the lab.}
\label{tab:ae-inputs}
\begin{tabular}{llrr}
\toprule
 & & minute 0 & minute 240 \\
\midrule
Concentrate & tare & \multicolumn{2}{c}{$96.2\unit{g}$} \\
            & weight            & $2119.7\unit{g}$ & $2101.5\unit{g}$ \\
            & lithium           & $63.56\unit{mg/L}$ & $218.5\unit{mg/L}$ \\
            & magnesium         & $2.974\unit{mg/L}$ & $130.0\unit{mg/L}$ \\
\addlinespace
Dilute      & tare & \multicolumn{2}{c}{$111.5\unit{g}$} \\
            & weight            & $2117.8\unit{g}$ & --- \\
            & lithium           & $403.1\unit{mg/L}$ & --- \\
            & magnesium         & $1657\unit{mg/L}$ & --- \\
\bottomrule
\end{tabular}
\end{table}

The dilute's last minute is left out of the table on purpose. Recovery divides by what the
feed held \emph{before} anything moved, so only its first sampled minute is needed; its last
minute is used by the balance gap, which is a diagnostic. The two tares matter as much as the
weights do, because neither tank weight means anything until its own empty weight is taken
off it.

\subsection{Weights become litres}

Nothing measured a volume. Each tank sat on a balance all run, so
\texttt{VolumeInATank}'s \texttt{find\_litres} turns a reading into litres by subtracting
that tank's tare and dividing by $1000\unit{g/L}$ --- the dilute salt solutions in these tanks weigh
close enough to a kilogram per litre that the balance can be read as a volume.

\begin{align}
V &= \frac{w - w_{\mathrm{tare}}}{1000} \label{eq:volume-ae} \\
V_{\mathrm{conc}}(0)   &= \frac{2119.7 - 96.2}{1000}  = 2.0235\unit{L} \nonumber \\
V_{\mathrm{conc}}(240) &= \frac{2101.5 - 96.2}{1000}  = 2.0053\unit{L} \nonumber \\
V_{\mathrm{dil}}(0)    &= \frac{2117.8 - 111.5}{1000} = 2.0063\unit{L} \nonumber
\end{align}

The concentrate lost $18.2\unit{g}$ over the run, about nine parts in a thousand of its
volume. That is why \eqref{eq:volume-ae} is run at each sampled minute rather than once: a
concentration compared minute to minute against a volume that has shifted is comparing two
different amounts of liquid.

\subsection{Litres and concentrations become milligrams}

\texttt{AmountInATank.find\_milligrams} multiplies the lab's $\unit{mg/L}$ by the litres from
\eqref{eq:volume-ae}. Four numbers come out of AE's concentrate:

\begin{align}
m &= c \, V \label{eq:amount-ae} \\
\mathrm{Li}(0)   &= 63.56 \times 2.0235 = 128.614\unit{mg} \nonumber \\
\mathrm{Li}(240) &= 218.5 \times 2.0053 = 438.158\unit{mg} \nonumber \\
\mathrm{Mg}(0)   &= 2.974 \times 2.0235 =   6.018\unit{mg} \nonumber \\
\mathrm{Mg}(240) &= 130.0 \times 2.0053 = 260.689\unit{mg} \nonumber
\end{align}

The feed's first minute goes the same way: $403.1 \times 2.0063 = 808.740\unit{mg}$ of
lithium and $1657 \times 2.0063 = 3324.439\unit{mg}$ of magnesium.

\subsection{What crossed}

\texttt{AmountThatCrossed} takes the sampled minutes of one tank, sorts them \emph{by
minute}, and subtracts the first from the last. It sorts because the readings arrive in file
order, which is not guaranteed to be time order; taking the first and last rows as they came
would silently reverse the sign on a run whose rows were entered out of sequence.

\begin{align}
\Delta m_{\mathrm{Li}} &= 438.158 - 128.614 = 309.544\unit{mg}
   &&\rightarrow\; \frac{309.544}{6.94} = 44.603\unit{mmol} \label{eq:crossed-li} \\
\Delta m_{\mathrm{Mg}} &= 260.689 -   6.018 = 254.671\unit{mg}
   &&\rightarrow\; \frac{254.671}{24.305} = 10.478\unit{mmol} \label{eq:crossed-mg}
\end{align}

Read in milligrams, magnesium looks almost as mobile as lithium. Read in millimoles it does
not: lithium crossed in four times the molar quantity. That is the whole reason the code
converts before it compares, and it is why the next two subsections look so different from
each other.

\subsection{Recovery}

\texttt{LithiumRecovery.find\_percent} divides what reached the concentrate by what the feed
started with. Both sides are milligrams, so the molar mass cancels and there is nothing to
gain from going through moles.

\begin{equation}\label{eq:recovery-ae}
R_{\mathrm{Li}} = 100 \times \frac{\Delta m_{\mathrm{Li}}}{m_{\mathrm{dil,Li}}(0)}
                = 100 \times \frac{309.544}{808.740} = 38.275\,\%
\end{equation}

The denominator is the feed's \emph{first} minute, never its last. Against the
$505.257\unit{mg}$ the feed ended with, the same crossing would read $61.3\,\%$, which would
reward a run for losing feed to the sample bottles.

\subsection{The charge that passed}

No meter totalled the charge. The supply current was written down every five to fifteen
minutes, so \texttt{ChargeThatPassed} integrates those readings with the trapezoid rule:
average the two currents across a gap and multiply by its length.

\begin{equation}\label{eq:charge-ae}
Q = 60 \sum_{k} \frac{I_k + I_{k+1}}{2}\,(t_{k+1} - t_k)
\end{equation}

The factor of 60 is the only unit conversion in \eqref{eq:charge-ae}: the log counts minutes and
a coulomb is an amp-second. AE's log holds 30 minutes with a current, from minute 0 to minute
240. The first gap contributes $\tfrac{1}{2}(0.081 + 0.082) \times 10 = 0.815\unit{A\,min}$,
the last $0.082 \times 15 = 1.23\unit{A\,min}$, and all 29 gaps together come to
$19.6425\unit{A\,min}$. So

\begin{align}
Q &= 19.6425 \times 60 = 1178.55\unit{C} \label{eq:ae-charge} \\
\frac{Q}{A} &= \frac{1178.55}{3600 \times 6.0} \times 1000 = 54.563\unit{mA\,h/cm^2}
   \nonumber
\end{align}

Minute 170 of AE's log has a voltage but no current, so it never enters \eqref{eq:charge-ae} and
the curve bridges minute 165 to minute 180 in one step. The run's current density is a
separate number: \texttt{CurrentDensity} takes the \emph{median} of the same currents,
$0.082\unit{A}$, giving $0.082 \times 1000 / 6.0 = 13.667\unit{mA/cm^2}$. The median rather
than the mean, because one run holds one current all the way through and a single misread
minute should not move the label on the whole run.

\subsection{The transport numbers}

A transport number is a share of the \emph{charge}, not of the ions. Every magnesium ion
carries two charges, so \texttt{TransportNumbers} weights each metal's millimoles by the
charge on its own ion. The charge sits on the \texttt{Metal} object next to the molar mass,
which is what stops the weighting from being forgotten at a call site.

\begin{align}
z_{\mathrm{Li}} n_{\mathrm{Li}} &= 1 \times \frac{309.544}{6.94}   = 44.603 \label{eq:tn} \\
z_{\mathrm{Mg}} n_{\mathrm{Mg}} &= 2 \times \frac{254.671}{24.305} = 20.956 \nonumber \\
t_{\mathrm{Li}} &= \frac{44.603}{44.603 + 20.956} = \frac{44.603}{65.559} = 0.6803
   \nonumber \\
t_{\mathrm{Mg}} &= \frac{20.956}{65.559} = 0.3197 \nonumber
\end{align}

Lithium crossed in $4.26$ times the molar quantity magnesium did, yet carried only $2.13$
times the charge. Doubling magnesium in \eqref{eq:tn} halves that advantage, and the two
shares add to one by construction.

\subsection{The feed ratio and permselectivity}

$t_{\mathrm{Li}}/t_{\mathrm{Mg}}$ on its own says as much about the feed as about the
membrane: a feed with little magnesium in it cannot send much magnesium across. So
\texttt{Permselectivity} divides the transport ratio by the ratio the feed actually offered,
which \texttt{FeedRatio}'s \texttt{find\_li\_over\_mg\_\allowbreak in\_the\_dilute} measures
from the dilute's first sampled minute. Both amounts come from one tank at one minute, so
the litres cancel and the milligrams carry the ratio.

\begin{align}
r_{\mathrm{feed}} &= \frac{808.740 / 6.94}{3324.439 / 24.305}
                   = \frac{116.533}{136.780} = 0.85197 \label{eq:feed-ratio} \\
P &= \frac{t_{\mathrm{Li}} / t_{\mathrm{Mg}}}{r_{\mathrm{feed}}}
   = \frac{0.6803 / 0.3197}{0.85197} = \frac{2.12838}{0.85197} = 2.4982 \label{eq:perm-ae}
\end{align}

Dividing by the feed ratio pins ``no preference'' at $0.5$ whatever the feed was, which is
the only thing that lets AE and AG be set beside each other at all. AE's $2.4982$ says the
stack favoured lithium about five times over what the feed handed it.

Note that \eqref{eq:feed-ratio} is lithium over magnesium, measured in the tank. The report's
\texttt{Mg/Li} column is a different number. \texttt{FeedRatio}'s other method,
\texttt{find\_mg\_over\_li\_\allowbreak from\_the\_setup}, reads the run sheet's own feed
concentrations, the other way up, giving
$(2000/24.305)/(400/6.94) = 1.4277$ for AE. The two are kept apart by name because they
answer different questions --- the measured one is what \eqref{eq:perm-ae} divides by, the set-up
one is a design column of the polynomial.

\subsection{The voltage and the energy}

AE's power log never passed $16.0\unit{V}$, reached at minute 0, and fell to
$11.31\unit{V}$ by minute 240 as the concentrate gained salt and its resistance dropped. The
highest voltage is what the $20\unit{V}$ supply ceiling is judged against, and on a healthy
run it is the starting voltage.

Energy uses the same trapezoid as \eqref{eq:charge-ae}, over the product of volts and amps, and
only over minutes that carry \emph{both}. AE has 30 such minutes, totalling
$245.996\unit{W\,min}$:

\begin{align}
E &= 245.996 \times 60 = 14759.78\unit{J} = 4.0999\unit{W\,h} \label{eq:energy-ae} \\
n_{\mathrm{Li}} &= \frac{309.544}{1000 \times 6.94} = 0.044603\unit{mol} \nonumber \\
\frac{E}{n_{\mathrm{Li}}} &= \frac{4.0999}{0.044603} = 91.92\unit{W\,h/mol} \nonumber
\end{align}

Per mole of lithium delivered, rather than per hour or per litre, so a short run and a long
one, or a gentle current and a hard one, can be compared on the same axis. The denominator is
the concentrate's gain from \eqref{eq:crossed-li}, and a run whose concentrate gained no
lithium gets \texttt{None} here instead of a division by zero.

\subsection{Run AG beside it}

AG is AE with one setting changed. Same GO-10 membrane, same $6.0\unit{cm^2}$, same
$400\unit{mg/L}$ lithium feed, same $0.082\unit{A}$ held for 240 minutes --- and
$8000\unit{mg/L}$ magnesium instead of $2000$, which is four times the magnesium on a molar
basis too, $5.7108$ against $1.4277$.

\begin{table}[htbp]
\centering
\caption{AE and AG, every stage. Both ran GO-10 at $13.667\unit{mA/cm^2}$ for 240 minutes;
only the feed magnesium differs.}
\label{tab:ae-ag}
\begin{tabular}{lrr}
\toprule
 & AE & AG \\
\midrule
Feed magnesium ($\unit{mg/L}$)          & $2000$ & $8000$ \\
Feed Mg/Li from the run sheet           & $1.4277$ & $5.7108$ \\
\addlinespace
Concentrate tare ($\unit{g}$)           & $96.2$ & $120.6$ \\
Dilute tare ($\unit{g}$)                & $111.5$ & $128.8$ \\
Concentrate volume, minute 0 ($\unit{L}$)   & $2.0235$ & $2.0167$ \\
Concentrate volume, minute 240 ($\unit{L}$) & $2.0053$ & $1.9847$ \\
Dilute volume, minute 0 ($\unit{L}$)        & $2.0063$ & $2.0312$ \\
\addlinespace
Concentrate Li, minute 0 ($\unit{mg}$)   & $128.614$ & $139.999$ \\
Concentrate Li, minute 240 ($\unit{mg}$) & $438.158$ & $193.488$ \\
Concentrate Mg, minute 0 ($\unit{mg}$)   & $6.018$ & $0.190$ \\
Concentrate Mg, minute 240 ($\unit{mg}$) & $260.689$ & $615.455$ \\
Dilute Li, minute 0 ($\unit{mg}$)        & $808.740$ & $868.135$ \\
Dilute Mg, minute 0 ($\unit{mg}$)        & $3324.439$ & $14129.027$ \\
\addlinespace
Li crossed ($\unit{mg}$ / $\unit{mmol}$) & $309.544$ / $44.603$ & $53.489$ / $7.707$ \\
Mg crossed ($\unit{mg}$ / $\unit{mmol}$) & $254.671$ / $10.478$ & $615.265$ / $25.314$ \\
\addlinespace
Lithium recovery                         & $38.275\,\%$ & $6.161\,\%$ \\
Charge passed ($\unit{C}$)               & $1178.55$ & $1179.90$ \\
Charge passed ($\unit{mA\,h/cm^2}$)      & $54.563$ & $54.625$ \\
$t_{\mathrm{Li}}$ / $t_{\mathrm{Mg}}$    & $0.6803$ / $0.3197$ & $0.1321$ / $0.8679$ \\
Feed Li/Mg in the dilute                 & $0.85197$ & $0.21518$ \\
Permselectivity                          & $2.4982$ & $0.70745$ \\
\addlinespace
Highest voltage ($\unit{V}$)             & $16.0$ & $13.9$ \\
Energy ($\unit{W\,h/mol}$ of Li)         & $91.92$ & $474.68$ \\
Li balance gap                           & $-19.645\,\%$ & $18.406\,\%$ \\
Mg balance gap                           & $35.002\,\%$ & $44.626\,\%$ \\
\bottomrule
\end{tabular}
\end{table}

The two runs pushed essentially the same charge --- $1178.55$ against
$1179.90\unit{C}$, a tenth of a percent apart --- and spent it on completely different work.
Lithium recovery fell from $38.3\,\%$ to $6.2\,\%$, because magnesium took $86.8\,\%$ of the
charge instead of $32.0\,\%$ and the lithium simply had no current left to ride on. The cost
of a mole of lithium rose from $92$ to $475\unit{W\,h}$ for the same reason: the energy
denominator in \eqref{eq:energy-ae} shrank by a factor of nearly six while the energy itself
fell only a tenth.

Permselectivity tells the more careful story. It fell too, from $2.4982$ to $0.70745$, but
far less than recovery did --- a factor of $3.5$ against a factor of $6.2$. Dividing by the
feed ratio in \eqref{eq:perm-ae} takes out most, though plainly not all, of the four-fold
magnesium increase. What is left is a real loss of selectivity under a magnesium-rich feed,
not an artefact of the feed being different, and that is exactly the separation
\eqref{eq:perm-ae} exists to make.

One detail of AG is worth following, because it shows the checks and the calculations
working together. Minute 225 of AG's power log reads $100.06\unit{V}$ on the supply against
$10.0\unit{V}$ on the multimeter. \texttt{PairedValuesTooFarApart} clears the supply reading
and writes it to the journal; had it not, AG's highest voltage would read $100.06\unit{V}$
and its energy would be overstated. With the reading gone, the energy trapezoid bridges
minute 210 to minute 240 in one $30$-minute step worth $24.76\unit{W\,min}$, and the highest
voltage falls back to the $13.9\unit{V}$ at minute 0, which is what the table reports.

The balance gaps are diagnostics only; nothing is ever corrected by them. AE's lithium gap is
negative, meaning the concentrate and rinse together ended with $59.6\unit{mg}$ more lithium
than the feed appears to have lost, which is sampling and weighing error rather than lithium
appearing from nowhere. The magnesium gaps of $35\,\%$ and $45\,\%$ are mostly liquid left
behind in the sample bottles, and the rinse is counted as zero whenever it was never weighed,
which can only make a gap look worse --- the honest direction for a diagnostic to be wrong in.

\subsection{What the test suite pins}

The numbers above are AE's and AG's, and they are not literals anywhere in the repository ---
re-run \texttt{uv run limodel analyze} and the report recomputes them from the four CSV files
every time. What the suite pins is the arithmetic that produced them, one test file per
quantity under \texttt{tests/\allowbreak analyzer/\allowbreak calculate/\allowbreak
quantities/}, built on fixture runs whose
tares are round and whose tanks hold exactly $1.000$ or $2.000\unit{L}$ so every expected
milligram can be checked in the head.

So each choice this section leaned on has a test standing behind it. A fixture holding
$10\unit{mmol}$ each of lithium and magnesium asserts $t_{\mathrm{Li}} = 1/3$, which fails the
moment the charge weighting in \eqref{eq:tn} is dropped. Another asserts recovery of
$30\,\%$ and states in its own name that it measures against the feed's first minute and not
its last, so moving the denominator in \eqref{eq:recovery-ae} breaks the build. A third feeds the
same crossing from two different feeds and asserts permselectivity of $0.5$ and then $1.0$,
which is \eqref{eq:perm-ae} refusing to be rewritten without the feed ratio. The trapezoid in
\eqref{eq:charge-ae} is pinned by a steady $0.1\unit{A}$ for an hour reading $360\unit{C}$, and
by a log handed over out of minute order that must still read $900\unit{C}$.

A change that quietly altered any step of AE's chain would have to get past those first.
